% Opciones empleadas:
%
%	a4paper -> indica el tama�o del papel, en este caso A4.
%	11pt -> tama�o de la fuente 11 puntos.
%	oneside -> s�lo escribimos en una cara del folio.
% 
% Otras opciones interesantes:
%
%	twoside -> escribimos a doble cara.
%	openbib -> para que las referencias bibliogr�ficas tengan un salto de l�nea entre cada campo de la referencia.
%
\documentclass[a4paper,11pt,twoside]{book}

% codificaci�n latin1 y s�mbolos del idioma espa�ol (�, acentos, ...)
\usepackage[english]{babel}
\usepackage[latin1]{inputenc}

% puede que queramos usar el s�mbolo del euro.
\usepackage{eurosym}

% El paquete fancybox nos permite crear cajas de diferentes estilos con facilidad.
% http://www.ctan.org/get/macros/latex/contrib/fancybox/fancybox.pdf
% http://www.mackichan.com/index.html?techtalk/487.htm~mainFrame
\usepackage{fancybox}
\usepackage{multicol}
\usepackage{amsmath}

% Para incluir subfiguras.
\usepackage{subfigure}

% Para incluir gr�ficos en JPG => compilar con pdflatex.
% \usepackage[pdftex]{graphicx}
\usepackage[pdftex]{graphicx}   % --> LaTex =>PDF con epstodf
\usepackage{epstopdf}

% Para incluir gr�ficos EPS => compilar con latex.
%\usepackage[dvips]{graphicx}

% Para escribir en color...
%
% ... cuando compilamos con el comando ``latex''
%\usepackage[dvips,usenames]{color}
% ... uando compilamos con el comando ``pdflatex''
\usepackage[pdftex,usenames,dvipsnames]{color}

% Espaciado y ajuste de m�rgenes
\usepackage{setspace}
\onehalfspacing
% \doublespacing
\setlength{\textwidth}{15cm}
\setlength{\textheight}{22cm}
\setlength{\oddsidemargin}{1cm}
\setlength{\evensidemargin}{1cm}
%\usepackage[hmargin={4cm,2.5cm},tmargin=4cm,bmargin=4cm]{geometry}
%Code snippets.
\usepackage{listings}
%Landscape figures.
\usepackage{lscape}

% Paquete fancyhdr -> Para modificar la cabecera y pie de p�ginas.
% http://tug.ctan.org/tex-archive/macros/latex/contrib/fancyhdr/
\usepackage{fancyhdr}
\pagestyle{fancy}
\fancyhf{}
%\fancyhf[HR]{\thepage}
%\fancyhf[HL]{\nouppercase\rightmark}
\fancyhead[LO]{\leftmark}
\fancyhead[RE]{\rightmark}
\fancyhead[RO,LE]{\thepage}
\renewcommand{\chaptermark}[1]{\markboth{\textbf{\thechapter. #1}}{}}
\renewcommand{\sectionmark}[1]{\markright{\textbf{\thesection. #1}}}
\renewcommand{\headrulewidth}{0.6pt}
\setlength{\headheight}{2.5\headheight}

% Package booktabs -> Para mejorar el aspecto de las tablas o cuadros.
% http://www.ctan.org/tex-archive/macros/latex/contrib/booktabs/
\usepackage{booktabs}

% Package rotating -> Para poder girar las tablas y dibujarlas a lo largo
% del folio en vez de a lo ancho.
\usepackage{rotating}

% Packages multicol y multirow, para manejar tablas de filas y columnas m�ltiples.
\usepackage{multicol}
\usepackage{multirow}
%To center the captions.
\usepackage[center]{caption}

%Links in biblio.
\usepackage{url}

\usepackage[pagebackref=true, hypertexnames=true, plainpages=false]{hyperref}

%Bold colors
\usepackage{color,soul}

%Math caligraphy
\DeclareMathAlphabet{\mathpzc}{OT1}{pzc}{m}{it}

%\usepackage[T1]{fontenc}

% Personalizamos la separaci�n entre p�rrafos...
\parskip=6pt

% Personalizamos el identado en la primera l�nea del nuevo p�rrafo...
\parindent=0pt

% Establecemos el n�mero m�ximo de niveles de profundidad en las secciones.
\setcounter{secnumdepth}{3}

% T�tulo
\title{Design and implementation of a HPC solution for stereo matching}
% Autor
\author{Luis Omar Alvarez Mures}
% Fecha
\date{\today}


\begin{document}

	% \maketitle sirve para generar autom�tica una portada predefinida, pero para un proyecto fin de carrera
	%	de FIC no sirvir�a porque no cumple las normas de presentaci�n. Podemos hacer dos cosas:
	% 1. Usarla e ignorar las normas (y asumir las consecuencias que pueda tener)
	% 2. Hacernos una portada en LaTeX que cumpla las normas (menos arriesgado)
	%
	
	%Modify list bullets so it does not use dash for second level
	\renewcommand{\labelitemi}{$\bullet$}
	\renewcommand{\labelitemii}{$\circ$}
	
	%Change page numbering so backref works
	\pagenumbering{roman}	
	
        %
% Portada.
%

% Nota: Ser�a m�s c�modo emplear el comando \maketitle que genera una portada de forma autom�tica, pero 
% no incluye toda la informaci�n que es necesario incluir en la memoria de un proyecto de fin de carrera
% de la Facultad de Inform�tica de A Coru�a.
%

\begin{titlepage}

  \begin{center}
    \includegraphics[width=7cm]{figures/anagramaUDC.pdf}\\[0.75cm]
    {\textsc{Faculty of Computer Science}} \\
    {\large \textsc{Departament of Electronics and Systems}} \\[1cm]
    {\Large \textsc{Final Year Project}} \\[0.25cm]
    {\Large \textsc{Master in High Performance Computing}} \\[2cm]
    {\LARGE \textsl{\textbf{Design and implementation of a HPC solution for stereo matching}}} 
    \vfill
    \begin{flushright}
      \begin{tabular}{ll}
        \textbf{Student:}    & Luis Omar �lvarez Mures\\
        \textbf{Advisors:} & Emilio J. Padr�n Gonz�lez\\
        & Juan Ram�n Rabu�al Dopico\\
        & \\
        \multicolumn{2}{r}{\small \emph{A Coru�a, \today{}.}} \\
      \end{tabular}
    \end{flushright}
  \end{center}

\end{titlepage}

        \thispagestyle{empty}
        \cleardoublepage
        %\newpage
	%\thispagestyle{empty}
	%\mbox{}

	% FRONTMATTER: TOC, LOF, LOT y descripci�n/organizaci�n de la memoria.
        \frontmatter
	
	% Los proyectos de fin de carrera de FIC han de ir acompa�ados de una serie de documentos adicionales, algunos
	% 	de ellos obligatorios (certificado, resumen, lista de palabras clave) y otros opcionales (dedicatoria
	%	y agradecimientos).
	%
	
        %\thispagestyle{empty}     % No number page, headings...
        %\include{tfg_certificado}
        
        \thispagestyle{empty}     % No number page, headings...
	%\include{tfg_desc}
	
	\thispagestyle{empty}
        \cleardoublepage
	
	\thispagestyle{empty}     % No number page, headings...
	%\include{tfg_dedicatoria}
        
        \thispagestyle{empty}
        \cleardoublepage
        
        %\newpage
	%\thispagestyle{empty}
	%\mbox{}
	
	\thispagestyle{empty}
        %\include{tfg_agradecimientos}
        
        \thispagestyle{empty}
        \cleardoublepage
        %\newpage
	%\thispagestyle{empty}
	%\mbox{}
	
	\thispagestyle{empty}
       % \include{tfg_resumo}
        
        \thispagestyle{empty}
        \cleardoublepage
	
	\thispagestyle{empty}
       % \include{tfg_resumen}
        
        \thispagestyle{empty}
        \cleardoublepage
        
        \thispagestyle{empty}     % No number page, headings...
	%\include{tfg_palabras_clave}
	
	\thispagestyle{empty}
        \cleardoublepage
	
        \thispagestyle{empty}     % No number page, headings...
       % \include{tfg_soft_hard}
        
        \thispagestyle{empty}
        \cleardoublepage
        
       % \thispagestyle{empty}     % No number page, headings...
        
        %\newpage
	%\thispagestyle{empty}
	%\mbox{}

        \tableofcontents
        \listoffigures
        \listoftables
        
        %\newpage
	%\thispagestyle{empty}
	%\mbox{}

	% MAINMATTER: El contenido, cap�tulo a cap�tulo, de la memoria del tfg.
        \mainmatter
	
	\pagenumbering{arabic}	
	
	%\include{tfg_frontmatter_intro}
	%\include{tfg_main_020}
	%\include{tfg_main_030}

	% INCLUIMOS LOS AP�NDICES...
        %\appendix
	%\include{tfg_appendix_010}
%         \include{tfg_appendix_020}

	%Primero hay que compilar el bibtex para que salga.
	% INCLUIMOS LA BIBLIOGRAF�A...
	\nocite{*}	% Se usa para indicar en la bibliograf�a las referencias no citadas.
	%\bibliography{tfg_biblio}
	\bibliographystyle{alpha}

\end{document}

